\section{Scope: what is established and how to read the output}\label{sec:scope}

\subsection{What rests on what}
\begin{itemize}
\item \emph{Proved.} For exact weights, known control coefficients and a fixed centring
constant: the variance estimate is exactly unbiased when movers' variances are known;
it converges under design conditions implied by KSS's; a wrong mover model adds the
explicit term $b(\Delta)$; and the same inputs give a consistent covariance matrix for
KSS's $q=1$ interval, with a coverage bound for the larger-of-two rule
(Appendices~\ref{sec:results} and~\ref{app:proofs}).
\item \emph{Shown by simulation.} The profile evaluated at estimated effects and the
projected weights (Section~\ref{sec:sim}); the $q=1$ interval at a bottleneck
(Section~\ref{sec:q1sim}); the automatic choice of interval and the control term
through the command (Section~\ref{sec:impl-mc}). These runs show calibration on the
designs that were run. They are not a coverage guarantee beyond them.
\item \emph{Compared.} Agreement with the KSS Matlab package and with KSS's
split-sample method on six real samples, and against a known truth on two of their
networks (Section~\ref{sec:kss}).
\end{itemize}

\subsection{Where the standard error was measurably off}
In most designs the standard error was within ten per cent of its benchmark and the
reported interval covered. The exceptions found are these.
\begin{itemize}
\item \emph{A few linking movers whose noise the model misses.} On the older
leave-two-out network, with wage noise tied to the weights that carry the component,
the normal-form standard errors for worker variance and the covariance were 14--17\%
too small. The reported $q=1$ intervals covered 96--98\%. In the synthetic designs with
noisy bridging workers the normal standard error was 0.59 to 0.67 times the true one,
and the $q=1$ interval covered 95.0--99.1\%. The exposure is to users who override the
default with \code{inference(highrank)} on a design with a large lead share.
\item \emph{Short panels.} In a three-year window the firm-variance standard error was
up to 16\% too large, beyond the registered 15\%. The error is on the conservative
side.
\item \emph{Estimated controls.} The standard error describes the fixed-offset
estimate. Where blocks of firms differ and a control loads on the worker effect, the
worker- and firm-variance standard errors were 7--8\% short in the registered
confirmation. With a control index as large as the effects, a stress case outside
that confirmation, the covariance standard error of the default joint estimate was
12\% short, and \code{nuisance(fixedoffset)} removed the shortfall. Variances that
move with a control inside a cell and poorly identified control coefficients are
further limits (Appendix~\ref{app:controls}).
\item \emph{A handful of bridging workers, or several weak links.} The coverage result
for the $q=1$ interval requires that no single row dominates the leading score. Two
groups of firms joined by three workers are outside it, although the interval covered
in the two such designs that were run. Two or more dominant directions are not
implemented.
\item \emph{Three registered checks failed and stand as failures:} the three-year
validation limit above, a length tolerance for the $q=1$ interval without a
bottleneck (12.7\% against 10\%), and the first package confirmation, on its
inventory.
\end{itemize}
Two things lie outside the standard error altogether. Independence of errors across
deletion units is the assumption of the point estimate itself; if it fails, the
estimate is affected and not only its precision. And the numerical noise of the point
estimate from random projections is reported separately: at the default 200
projections it was 10--12\% of the standard error on the older-worker samples, where
more projections are advisable.

\subsection{Reading the output}
\begin{itemize}
\item \emph{A normal interval with a small lead share} is the usual case on large,
well-connected panels. The standard error and \code{e(V)} can be used as usual, for
example with \code{test} and \code{lincom}.
\item \emph{A $q=1$ interval} means that one direction of the mobility network carries
more than a tenth of the estimated sampling variance, or that the normal variance
estimate was not positive. Report the interval. It has no standard
error, and \code{e(V)} is not posted.
\item \emph{With controls,} the standard error covers their estimation to first order.
If the control index is large relative to the effects, \code{nuisance(fixedoffset)}
under match deletion gives the estimate that the standard error describes.
\item \emph{The line on numerical precision} under the table says how large the
projection noise of the point estimate is relative to the standard error. Above about
a tenth, raise \code{probes()}.
\item \emph{\code{estat inference}} shows both intervals of every component, the lead
share, the two largest eigenvalues and the effective rank, and the warnings that
apply.
\end{itemize}
